%=============================================================================!
% 4.2  ENERGY IN THE OSCILLATOR                                               !
%=============================================================================!
\section{Energy in the oscillator}
\label{sec:os-energy}

An oscillating block has its largest kinetic energy at the centre of the
swing and none at the ends, and its potential energy does the opposite.
This section follows the trade and shows that, averaged over a period,
the energy is shared equally between the two forms, a result that
returns when Chapter~12 relates kinetic energy to temperature.

\subsection*{The trade}

With $x = A\cos(\omega t + \phi)$ and $v = -A\omega\sin(\omega t + \phi)$
(\cref{sec:os-circle}), and writing $\theta = \omega t + \phi$ for short,
the two energies of \cref{sec:en-potential} are
\begin{align*}
   K &= \tfrac12 m v^2 = \tfrac12 m A^2\omega^2\sin^2\theta
      = \tfrac12 k A^2\sin^2\theta
      && \text{since $\omega^2 = k/m$ gives $m\omega^2 = k$,}\\
   U &= \tfrac12 k x^2 = \tfrac12 k A^2\cos^2\theta .
\end{align*}
Taking out the common factor $\frac12 kA^2$ and using $\sin^2\theta +
\cos^2\theta = 1$, their sum is $E = \frac12 kA^2(\sin^2\theta +
\cos^2\theta) = \frac12 kA^2$, constant, as \cref{prop:en-energy-1d} requires. The energy
therefore fixes the amplitude, $A = \sqrt{2E/k}$: an oscillation with four
times the energy swings twice as far. For the same block
released instead from rest \SI{4}{\centi\metre} out, as in
\cref{fig:ne-spring}, $E =
\frac12\times50\times0.04^2 = \SI{0.04}{\joule}$. At $x =
\SI{2}{\centi\metre}$ the potential energy is $\frac12\times50\times0.02^2
= \SI{0.01}{\joule}$, and the remaining \SI{0.03}{\joule} is kinetic
(\cref{fig:os-energy}b). As the block swings, $K$ and $U$ rise and fall
twice in every period. By \cref{eq:mo-addition} with $\cos\pi = -1$ and
$\sin\pi = 0$, $\sin(\theta + \pi) = -\sin\theta$ and $\cos(\theta + \pi)
= -\cos\theta$, and squaring removes the minus signs, so $\sin^2\theta$
and $\cos^2\theta$ repeat when $\theta$ grows by $\pi$, half a turn
(\cref{fig:os-energy}a).

\begin{figure}[htb]
\centering
\includegraphics{ch04_oscillations/energy}
\caption{Energy in the oscillator: the block of \cref{sec:os-circle},
released instead from rest \SI{4}{\centi\metre} out, so that $E =
\SI{0.04}{\joule}$. (a) Kinetic energy $K$, potential energy $U$ and
their sum against time over two periods; the dotted line is $E/2$, the
average of each. (b) The energy diagram, $U = \frac12 kx^2$ with the line
$E$; at $x = \SI{2}{\centi\metre}$ the energy splits into $U =
\SI{0.01}{\joule}$ and $K = \SI{0.03}{\joule}$, and the dots mark the
turning points.}
\label{fig:os-energy}
\end{figure}

\Needspace{16\baselineskip}
\subsection*{Equal shares on average}

\begin{toolbox}[The average of a function]
The average of $n$ numbers is their sum divided by $n$. To average a
function $f(t)$ over an interval from $0$ to $t_1$, sample it at $n$ equally
spaced times $t_k$, a step $\tau = t_1/n$ apart, and average the
samples:
\begin{equation*}
   \frac1n\sum_{k} f(t_k) = \frac{1}{t_1}\sum_k f(t_k)\,\tau
   \;\longrightarrow\; \frac{1}{t_1}\int_0^{t_1} f(t)\,\dd t
   \qquad\text{as } n \to \infty ,
\end{equation*}
the first step multiplying and dividing by $\tau$ and using $n\tau = t_1$, the second the definition of the integral as a limit of such
sums (\cref{sec:mo-integrating}). The average of $f$ over the interval is
therefore $\frac{1}{t_1}\int_0^{t_1} f\,\dd t$, the area under its graph divided by
the length of the interval. We write it with an overbar, $\overline{f}$.
\end{toolbox}

To average $\cos^2\theta$ over a period, rewrite it. The addition formula
of \cref{sec:mo-trig} with both angles equal to $\theta$ gives
$\cos2\theta = \cos^2\theta - \sin^2\theta$, and
\begin{align*}
   \cos 2\theta &= \cos^2\theta - (1 - \cos^2\theta) = 2\cos^2\theta - 1
      && \text{replacing $\sin^2\theta$,}\\
   \cos^2\theta &= \tfrac12 + \tfrac12\cos 2\theta
      && \text{adding 1 to both sides and halving.}
\end{align*}
The second term averages to zero over a period, $t_1 = 2\pi/\omega$. With
$\theta = \omega t + \phi$, so that $2\theta = 2\omega t + 2\phi$, the
function $\sin(2\omega t + 2\phi)/(2\omega)$ has derivative $\cos(2\omega
t + 2\phi)$ by the chain rule, so the fundamental theorem of calculus
(\cref{sec:mo-integrating}) turns the integral into the change in that
function,
\begin{equation*}
   \int_0^{t_1} \cos(2\omega t + 2\phi)\,\dd t
   = \frac{\sin(2\omega t_1 + 2\phi) - \sin 2\phi}{2\omega} = 0 ,
\end{equation*}
because $2\omega t_1 = 4\pi$, two whole turns, after which the sine takes
its old value. Averaging the two terms of $\cos^2\theta$ separately,
\begin{align*}
   \overline{\cos^2\theta}
   &= \frac{1}{t_1}\int_0^{t_1}\tfrac12\,\dd t
      + \frac{1}{t_1}\int_0^{t_1}\tfrac12\cos(2\omega t + 2\phi)\,\dd t
      && \text{the integral of a sum,}\\
   &= \frac{1}{t_1}\times\frac{t_1}{2} + 0 = \frac12
      && \text{a rectangle, and the integral above,}
\end{align*}
and, since $\sin^2\theta = 1 - \cos^2\theta$, the same steps give
$\overline{\sin^2\theta} = 1 - \frac12 = \frac12$. The trade above gave
$K = E\sin^2\theta$ and $U = E\cos^2\theta$, with $E$ a constant that
comes outside each integral, so
\begin{equation}
   \overline{K} = \overline{U} = \tfrac12 E :
   \label{eq:os-equal-shares}
\end{equation}
over a period, a harmonic oscillator holds half its energy as kinetic
energy and half as potential energy. For the block, each averages
\SI{0.02}{\joule}, the dotted line of \cref{fig:os-energy}(a).

\notebook{04}{the spring with an animation, energy bars and its phase
portrait}

\begin{selfcheck}
At which positions does the block of \cref{fig:os-energy} have equal
kinetic and potential energy?
\end{selfcheck}
